% HD-S40. GPT-6 Astra Pro, /root/sixfolds; inherited effort not exposed.
% 21 November 2026. The actual normal closure and its characteristic diagonal kernel.
\subsection{HD-S40: the normal closure of bireflections}
\label{sec:bireflection-normal-closure}

Retain the actual finite Gaussian coefficient group, primitive
Jacobian corner and rank-nine period-family conventions of
HD-S36--HD-S38. For every coefficient rank \(m\le14\),
every such source has ordinary Hermitian determinant \([-1]\)
or \([-3]\). Paragraph5 proves the corresponding split
determinant conclusion for the seven other fields listed there.

For \(10\le m\le14\), the proof uses the intrinsic normal
subgroup
\[
 N=\langle x\in G:\operatorname{rank}(\rho(x)-I)=2\rangle.
                                                               \tag{S40.1}
\]
Outside the determinant conclusions already proved using proper
normalized algebras, it is absolutely irreducible. Paragraphs1--3
reduce Guralnick--Saxl's alternatives to \(r=m/2\) complex
blocks of dimension two, with block-permutation quotient \(S_r\).
Paragraph4 eliminates that last case by the actual characteristic
block kernel proved in HD-G28. The proof uses the actual normal
subgroup and its characteristic kernels, without assuming that
a chosen complex block system is defined over \(K=\mathbb Q(i)\).

\paragraph{1. The actual normal subgroup and its coefficient algebra.}
If \(G\) contains a reflection, HD-S36 already gives the
asserted determinant conclusion. Otherwise every inertia
codimension is at least two. The exact parabolic dimension
formula and the Hurwitz parameter bound give
\[
             \sum_j c_j=2m+6,\qquad b\ge12.                 \tag{S40.2}
\]
For \(m\le14\), all \(c_j\ge3\) would contradict
\(3b\ge36>2m+6\). Hence (S40.1) is nontrivial.
Conjugation preserves the rank of \(x-I\), so \(N\) is a
literal normal subgroup of \(G\).

The matrix span
\[
                 \mathcal A=K[N]\subset\operatorname{End}_K(V)
                                                               \tag{S40.3}
\]
is a unital semisimple algebra: it is the image of the finite
group algebra, whose invariant complements are obtained by
averaging a linear projection over \(N\). It is normalized by
\(G\). Since a nonidentity scalar in dimension \(m>2\)
cannot have codimension two, \(\mathcal A\ne K\).
If it is proper, HD-S38 gives the determinant conclusion with
its actual block-quotient and polarized-isogeny maps. Otherwise
\(\mathcal A=\operatorname{End}_K(V)\). Its complexification
is the full matrix algebra, so a complex \(N\)-invariant
subspace would be invariant under every matrix. Consequently
\(N\) is absolutely irreducible.

Guralnick--Saxl, \emph{Generation of finite almost simple groups
by conjugates}, J.~Algebra \textbf{268} (2003), 519--571,
Theorem8.4, printed pp.557--558, applies to this finite,
irreducible, bireflection-generated subgroup. In ranks greater
than eight it gives three alternatives. The first is a primitive
group with no fixed tensor decomposition; the second is a line
system with permutation quotient \(A_m\) or
\(S_m\); or two-dimensional blocks with permutation quotient
\(S_{m/2}\). The tensor-decomposable alternative in that
theorem has dimension four and does not occur here.

The first alternative would also make \(G\) primitive and
tensor-indecomposable. Indeed a \(G\)-block decomposition is
preserved by its subgroup \(N\), and \(N\)'s irreducibility
makes the action on its blocks transitive. A tensor decomposition
preserved by \(G\) is preserved by \(N\) through the literal
restriction of its tensor-factor action. These contradict the
respective properties of \(N\). HD-S37 excludes a primitive,
tensor-indecomposable \(G\) in these ranks. It remains to
analyze the line alternative without any descent assumption.

\paragraph{2. A scalar diagonal kernel is impossible.}
We prove a slightly wider elementary assertion. Let \(M\)
be a finite irreducible complex monomial group of degree
\(m\ge7\), generated by actual codimension-two elements.
Suppose its permutation quotient is \(A_m\) or \(S_m\),
and let \(T\) be its diagonal kernel. Then \(T\) is not
scalar.

Suppose, for a contradiction, that \(T\) is scalar.
For a monomial matrix, the fixed dimension on a permutation
cycle is either one or zero: writing the cycle maps in order,
the last equation says that their product fixes the first
coordinate. Therefore the permutation of a codimension-two
element \(x\) is a transposition, a three-cycle, or a product
of two disjoint transpositions. The identity permutation would
make \(x\) scalar, which is impossible.

Consider a transposition, say \((1\,2)\). Since \(x^2\)
is scalar and \(x\) fixes at least one unchanged coordinate,
\(x^2=I\). Its two-cycle has holonomy one and contributes
one fixed vector. Exactly one of the other \(m-2\) lines
therefore has eigenvalue \(-1\), and all the rest have
eigenvalue one. In particular,
\[
                         \operatorname{tr}(x)=m-4\ne0.      \tag{S40.4}
\]
Every lift of an element of the permutation centralizer conjugates
\(x\) to \(z x\) for a scalar \(z\in T\). Trace forces
\(z=1\). But that centralizer acts transitively on the
\(m-2\) unchanged lines, whereas its lifts would have to
preserve their unique negative line. This is impossible.
Here a transposition can occur only when the permutation quotient
is \(S_m\); its centralizer contains the entire symmetric
group on the unchanged points.

For a three-cycle or a double transposition, attaining
codimension two forces holonomy one on each moving cycle and
the identity on every unchanged line. All these permutations
are even. Since they generate the permutation quotient, that
quotient is \(A_m\).

Let \(H\) be the stabilizer of the first line, and let
\(\chi:H\to\mathbb C^*\) be its actual action on that
line. The map \(H\to A_{m-1}\) is onto. Since
\(A_{m-1}\) is perfect, \(H'\) maps onto it. On the
other hand, \(H'\subset\ker\chi\), and \(\chi\) is
faithful on the scalar subgroup \(T\). Thus
\[
              H'\simeq A_{m-1},\qquad H=T H'.              \tag{S40.5}
\]
The group \(H'\) acts trivially on the first line and
transitively on the other \(m-1\) lines. The stabilizer of
one of those lines is \(A_{m-2}\), which is perfect for
\(m\ge7\). Its one-dimensional character is therefore
trivial. Choose a nonzero vector on the second line and
transport it by \(H'\) to the other lines. Triviality of
the stabilizer character makes this independent of the chosen
transport. In this explicit basis,
\[
              H'=\{P_\pi:\pi\in A_{\{2,\ldots,m\}}\}.
                                                               \tag{S40.6}
\]
The first basis vector is still free to be rescaled.

Some generating bireflection moves the first line. If it has
a three-cycle permutation, use that element. If it is a double
transposition, relabel its support as
\((1\,2)(3\,4)\) and conjugate the actual element
\(P_{(2\,5\,6)}\in H'\) by it. The result is an actual
bireflection with a three-cycle permutation containing the first
point. Conjugating within \(H'\) and taking the inverse if
necessary, we obtain an element \(g\) with
\[
 ge_1=a e_2,\qquad ge_2=b e_3,\qquad
 ge_3=c e_1,\qquad abc=1,\qquad ge_j=e_j\ (j\ge4).
                                                               \tag{S40.7}
\]
Take \(h=P_{(2\,4\,5)}\in H'\). The element
\(g h g^{-1}\) has permutation \((3\,4\,5)\) and fixes
\(e_1\) exactly. By (S40.5) it lies in \(H'\), not just
in a scalar multiple of \(H'\). In the specified basis its
three nonzero cycle entries are
\[
       e_3\longmapsto b^{-1}e_4,\qquad
       e_4\longmapsto e_5,\qquad e_5\longmapsto b e_3.
                                                               \tag{S40.8}
\]
Comparison with (S40.6) gives \(b=1\). Replacing
\(e_1\) by \(a^{-1}e_1\) makes \(g\) the literal
permutation \((1\,2\,3)\). The subgroup generated by it
and (S40.6) is the full permutation \(A_m\): conjugation
gives every three-cycle containing the first point, and the
point stabilizer gives all the others.

Every generating bireflection of \(M\) now differs from its
literal permutation by a scalar. If that scalar were not one,
none of its at least \(m-4\) unchanged lines could contribute
a fixed vector; the moving cycles contribute at most two.
This contradicts fixed dimension \(m-2>2\). All generating
bireflections therefore belong to this permutation \(A_m\).
It fixes \(e_1+\cdots+e_m\), contradicting irreducibility.
This proves the assertion.

\paragraph{3. The nonscalar diagonal subgroup is intrinsic.}
Return to the line alternative for \(N\). Its diagonal
kernel \(T\) is abelian and normal, and is nonscalar by
paragraph2. Every abelian normal subgroup of \(N\) has
trivial image in \(A_m\) or \(S_m\), because these groups
have no nontrivial abelian normal subgroup for \(m\ge7\).
Hence every such subgroup is contained in \(T\). It follows
that \(T\) is the unique maximal abelian normal subgroup
of \(N\), and is characteristic in \(N\). Since \(N\)
is normal in \(G\), the same actual \(T\subset GL_m(K)\)
is normal in \(G\).

Consequently
\[
                  K\subsetneq K[T]\subsetneq
                        \operatorname{End}_K(V)             \tag{S40.9}
\]
is a literal normalized semisimple algebra. The first inclusion
is strict because \(T\) is nonscalar; the second because
\(K[T]\) is commutative and \(m>1\). Its construction
does not require the individual complex lines or their
projection operators to be Galois stable. HD-S38 applies,
with the precise primitive-corner and polarized quotient maps
already proved there. This eliminates the line alternative and
proves the stated reduction, in particular the complete
coefficient-rank-11 and coefficient-rank-13 conclusion.

\paragraph{4. The characteristic two-block kernel closes the even ranks.}
The remaining GS alternative has \(m=2r\) with
\(r=5,6,7\). Apply the complete group-theoretic result
HD-G28 to the actual \(N\subset GL_m(K)\).
We recall its literal maps to specify the algebra being used.
Every outside bireflection exchanges exactly two blocks, is
the identity elsewhere, and its two transition maps are inverse.
Along a spanning tree of those exchanges, recursively set
\(b_j=x_{ij}|_{V_i}b_i\). The resulting actual complex
isomorphism \((u_j)\mapsto\sum b_j(u_j)\) makes all tree
edges pure swaps and embeds a pure \(S_r\subset N\).
Writing \(T\) for the block kernel gives \(N=T\rtimes S_r\).

Replacing any outside generator by its product with the inverse
pure swap, and retaining inside generators, shows that \(T\)
is generated by elements supported on at most two blocks.
If \(L\subset GL_2(\mathbb C)\) is their common coordinate
image, two elements whose other support blocks are distinct
give the exact commutator
\[
 [\operatorname{diag}(a,c,1,\ldots),
   \operatorname{diag}(b,1,d,\ldots)]
          =\operatorname{diag}([a,b],1,1,\ldots).
                                                               \tag{S40.10}
\]
Normal closure in each coordinate proves \((L')^r\subset T\).
The group \(L\) is irreducible: a common invariant line
would give an \(N\)-invariant sum of that line in every
block.

If \(L\) is solvable, \(T=\operatorname{Rad}(N)\), since
\(S_r\) has no nontrivial solvable normal subgroup.
Otherwise the actual projective map gives
\(L/(L\cap\mathbb C^*)=A_5\), by the finite \(PGL_2\)
classification cited and checked in HD-G28.
With \(Z=L\cap\mathbb C^*\) and \(Z_0=T\cap Z^r\),
the coordinatewise projective map \(\pi\) gives the literal
isomorphism
\[
 N/Z_0\longrightarrow A_5^r\rtimes S_r,\qquad
                    (tp)Z_0\longmapsto(\pi(t),p).            \tag{S40.11}
\]
Its surjectivity follows from \((L')^r\subset T\), and
its kernel is exactly \(Z_0\). The latter is
\(\operatorname{Rad}(N)\). The base \(A_5^r\) is the
unique minimal normal subgroup of this quotient: its factors
are transitively permuted and its centralizer is trivial.
Thus \(T\) is the inverse image of that intrinsic subgroup
and is characteristic in \(N\), including \(r=5\).

In either case the \emph{original} matrices of \(T\) span
a proper nonscalar semisimple \(K\)-algebra normalized by
\(G\). Properness follows after complexification from its
block-diagonal action; nonscalarity follows from irreducibility
of \(L\). No \(K\)-definition of the auxiliary tree basis
is needed. HD-S38 therefore gives \([-1]\) or \([-3]\),
eliminating the last case. HD-S36 and HD-S37 cover \(m\le9\),
so the Gaussian conclusion holds for every \(m\le14\).

\paragraph{5. The seven remaining quadratic fields and the exact geometry.}
The scalar-kernel argument in paragraph2 is field-free. The
same normal-subgroup proof therefore applies to
\[
 K=\mathbb Q(\sqrt{-d}),\qquad
                 d\in\{2,7,11,19,43,67,163\}.              \tag{S40.12}
\]
For these fields, every specified source with \(m\le14\)
has ordinary determinant \([-1]\), hence is split by
HD-S03a. In fact no source of this specified kind occurs in
these ranks for \(d\in\{11,19,43,67,163\}\).
Here the source conventions are the same actual full finite
group sphere covers, exact coefficient character field,
primitive group-algebra Jacobian corner, restricted theta form,
signature \((3,3)\), and period rank nine.

We give the precise replacement for the Gaussian proper-algebra
step. For \(10\le m\le14\), HD-S39 excludes a proper
nontrivial normalized semisimple coefficient algebra when the
reflection-free \(G\) is \(K\)-primitive. Otherwise there
is an actual transitive \(K\)-block system. A block \(U\)
has rank \(r\le7\). Its stabilizer \(H\) acts absolutely
irreducibly on it, and its character field is exactly \(K\)
by the literal induced trace formula. Put
\(H_0=\ker(H\to GL(U))\). The actual curves and coefficient
maps are
\[
 \begin{gathered}
 Y=\widetilde C/H,\qquad C'=\widetilde C/H_0,\\
 F_v(x)=\pi_U\rho(x)v,\qquad
 F\longmapsto\sum_{x\in H\backslash G}\rho(x)^{-1}F(x).
 \end{gathered}                                               \tag{S40.13}
\]
The inversion label \(Ht_j^{-1}\leftrightarrow t_jH\)
identifies the quotient sheets with the induced blocks.
The original invariant coefficient form makes the blocks
orthogonal, by averaging a positive block-diagonal form and
using absolute irreducibility. For \(v\in U(K)\setminus0\),
the normalized rational matrix-coefficient idempotents therefore
satisfy \(e_G=e_H\): coefficients off \(H\) vanish and
\(m/|G|=r/|H|\). The actual quotient pullback restricts to
a \(K\)-linear isogeny on the primitive factors and obeys
\[
 \alpha_*\alpha^*=[|H_0|],\qquad
 (\alpha^*)^*\theta_{J\widetilde C}=|H_0|\theta_{JC'},\qquad
          |H_0|^6=N_{K/\mathbb Q}(|H_0|^3).                 \tag{S40.14}
\]
The \(K\)-corner and action conventions are precisely the
class-sum and coefficient maps of HD-G27. These identities,
after fixing the finite labels and lattices, identify the
period map with the block-factor period map composed with the
actual parameter map.

Rank one is impossible, since \(\mu(K)=\{\pm1\}\) would
make the induced character rational. Rank two has period rank
at most seven. For ranks three through seven the exact formula
\(6=r(2g(Y)-2)+\sum c'_j\) and the marked-curve dimension
bound force \(g(Y)=0\) and \(b'\ge12\). Since
\(2r+6\le20<24\), at least one local monodromy is a
reflection. HD-G27 then gives determinant \([-1]\), or
excludes that faithful block entirely for the five stated
fields. The multiplier in (S40.14) preserves that determinant
class. This proves the required proper-algebra implication over
each field in (S40.12).

The primitive, tensor-indecomposable alternative for \(N\)
is excluded for these fields by the literal restriction argument
of paragraph1 and HD-S39. The line alternative gives the same
characteristic \(T\) and proper \(K[T]\) by paragraphs2--3.
The two-block alternative gives it by paragraph4 and HD-G28,
whose group-theoretic proof holds over every characteristic-zero
coefficient field. Thus the preceding proper-algebra implication
closes all ranks \(10\le m\le14\). HD-G27 treats
\(m\le8\), and HD-S39 excludes the reflection-free rank-nine
case; its reflection-containing case is again HD-G27.
This proves the asserted complete scope through rank fourteen.

\paragraph{Checks and source scope.}
The elementary scalar-kernel argument above does not use a
Schur-multiplier classification or an unproved descent of
imprimitivity systems. The new exact certificate checks the
weighted three-cycle conjugation, its basis rescaling, the
double-transposition conversion, and the finite rank budget.
The complete argument, rather than the finite calculations,
proves the universal assertions. The author independently read
the entire HD-G28 proof and its exact coordinate and quotient
maps, reopened its primary finite-\(PGL_2\) locator, and read
and freshly reran its supplementary certificate. The new
all-field conclusion uses HD-S39 and HD-G27 only through the
literal algebra, induced block, quotient-curve and polarized
isogeny maps stated above. No claim is made here about arbitrary
reflection-free coefficient ranks at least fifteen.
